\section*{Заключение}

В условиях цифровизации финансового сектора и роста долговой нагрузки населения, коммерческим банкам необходимо более точно и быстро оценивать кредитоспособность заемщиков. Традиционный подход к их оценке не способен выявлять сложные нелинейные зависимости в характеристиках клиента, а современный подход с применением моделей МО и ГО, наоборот, позволяет их выявить и повысить качество прогнозирования дефолта. 

Для достижения поставленной цели было проделано следующее:
\begin{enumerate}[nosep]
	\item в главе 1 изучены теоретические подходы к кредитному скорингу, дано определение на основе анализа научной литературы, рассмотрены статьи закона «Об искусственном интеллекте», косвенно затрагивающих применение ИИ в кредитном скоринге;
	\item в главе 2 проведен анализ исходных обезличенных данных кредитных историй АО «Альфа-Банка», проведены предобработка данных и сокращение признакового пространства, введено и обосновано определение аномального клиента и далее протестированы модели МО и ГО на 1 миллионе клиентов.
	\item в главе 3 данные модели были протестированы на больших данных и проведено сравнение оченки качества моделей МО и ГО.
\end{enumerate} 

В процессе исследования было установлено, что прогностическая способность моделей в большей степени зависит именно от качества подготовки данных, в отличие от подбора использования самих моделей. Значимым результатом данной работы является введение понятия аномального клиента как заемщика с фактической меткой недефолтного (платежеспособного) клиента, но при этом его платежная дисциплина схожа с платежным поведением дефолтных (неплатежеспособных) клиентов. Гипотеза о том, что аномальные клиенты по своему поведению схожи с дефолтными клиентами была подтверждена в результате статистической проверки теста Колмогорова-Смирнова. Проведенная переклассификация позволила качественно подготовить данные заемщиков для последующего моделирования. Все рассмотренные алгоритмы продемонстрировали высокие метрики точности.

Результаты показали, что наилучшими алгоритмами в данном исследовании являются деревья решений и случайный лес. Эти модели показали меньшее количество ошибочных предсказаний по сравнению с градиентным бустингом и нейронной сетью. При этом следует отметить, что дополнительный подбор гиперпараметров для градиентного бустинга не потребовался, поскольку изначально данный алгоритм не выявил эффекта переобучения по сравнению с деревьями решений и случайным лесом. Поэтому в данном случае алгоритм градиентного бустинга является наиболее устойчивым

Также метрики качества моделей по выявлению дефолтных клиентов оказались немного ниже, чем для недефолтных. Отсюда следует вывод, что распознование неплатежеспособных заемщиков является более трудоемкой задачей по сравнению с выявлением платежеспособных. 



